\documentclass[10pt]{article}

\usepackage[utf8]{inputenc}

\usepackage[T1]{fontenc}

\usepackage{amsmath}

\usepackage{amsfonts}

\usepackage{amssymb}

\usepackage[version=4]{mhchem}

\usepackage{stmaryrd}



\begin{document}

При целом $m$ справедлива формула



\[

\int_{0}^{1} e^{2 \pi i \alpha m} d \alpha= \begin{cases}1, & \text { если } m=0, \\ 0, & \text { если } m \neq 0,\end{cases}

\]



поэтому, если через $I(N)$ обозначить число решений уравнения



\[

N=u_{1}+\ldots+u_{k}, u_{v} \in U_{v}, v=1, \ldots, k,

\]



где $U_{v}$ - нек-рые подмножества натурального ряда чисел, то



\[

I(N)=\int_{0}^{1} S_{1}(\alpha) \ldots S_{k}(\alpha) e^{-2 \pi i \alpha N} d \alpha

\]



где



\[

S_{v}(\alpha)=\sum_{u_{v} \in U_{v}} e^{2 \pi i \alpha u_{v}}, v=1, \ldots, k

\]



Полагая, в тастности, $U_{v}=U=\left\{1^{n}, 2^{n}, \ldots\right\}$, получают Варинга проблему; при $k=3, U_{v}=U=\{2,3,5,7$, $11, \ldots, p, \ldots\}$ - тернарнуіо Гольдбаха проблему и т. Д. Как и в задаче о распределении дробных долей, главным вопросом здесь является вопрос об оценке сверху модуля $S_{v}(\alpha)$, т. е. вопрос о верхней грани $\left|S_{v}(\alpha)\right|$



Тем самым широкий круг разнообразных проблем теории чисел формулируется единообразно на языке тригонометрич. сумм.



Первая нетривиальная тригонометрич. сумма появилась у K. Гаусcа (C. Gauss, 1811) в одном из его доказательств закона взаимности квадратичных вычетов:



\[

S=\sum_{x=1}^{P} e^{2 \pi i F(x)}, F(x)=\frac{a x^{8}}{P}

\]



и К. Гаусс точно вычислил значение $S$ :



\[

S=\left\{\begin{array}{r}

(1+i) \sqrt{P}, \text { если } P \equiv 0 \bmod (4) \\

\sqrt{P}, \text { если } P \equiv 1(\bmod 4) \\

0, \text { если } P \equiv 2(\bmod 4) \\

i \sqrt{P}, \text { если } P \equiv 3(\bmod 4)

\end{array}\right.

\]



Целый ряд независимых работ с применением тригонометрич. сумм появился в нач. 20 в.



Г. Вейль (H. Weyl, 1916), изучая распределение дробных долей многочлена



\[

f(x)=\alpha_{1} x+\ldots+\alpha_{n} x^{n}

\]



с действительными коэффициентами $\alpha_{1}, \ldots, \alpha_{n}$, рассмотрел суммы $S$ с функцией $F(x)$ вида $F(x)=m f(x)$, $m$ - целое число, $m \neq 0$. И. М. Виноградов (1917) при изучении распределения целых точек в областях на плоскости и в пространстве рассмотрел суммы $S \mathrm{c}$ функцией $F(x)$, от к-рой требовалось только, чтобы ее вторая производная удовлетворяла условиям



\[

\frac{c_{1}}{A} \leqslant F^{\prime \prime}(x) \leqslant \frac{c_{2}}{A},

\]



где $c_{1}, c_{2}$ - нек-рые абсолютные положительные постоянные, $A=A(P) \rightarrow+\infty$ при $P \rightarrow+\infty$. Г. Харди (G. Hardy) и Дж. Литлвуд (J. Litlewood) (1918), получив приближенное функциональное уравнение дзета-функции Римана, рассмотрелии сумму $S$ с функцией $F(x)$ вида



\[

F(x)=t \ln x,

\]



где $t$ - действительный параметр, $t=t(P)$.



Во всех указанных работах требовалось найти возможно лучшую оценку сверху модуля суммы $S$.



Общая схема исследования названных проблем теории чисел Т. с. м. такова: вытисывается точная формула, выражающая число решений изучаемого уравнения или число дробных долей изучаемой функции, попадающих на заданный интервал, или число целых точек в заданной области, в виде интеграла от тригоно- метрич. суммы или в виде ряда, коәффициентами к-рого являются тригонометрич. суммы; точная формула представляется суммою двух слагаемых - главного и дополнительного (напр., если рассматривается ряд Фурье характеристич. функции интервала, то главный член получается от нулевого коәффициента ряда Фурье); главное слагаемое доставляет главный член асимптотич. формулы, дополнительное - остаточный член. В аддитивных задачах, таких, как проблема Варинга, проблема Гольдбаха и др., главное слагаемое исследуется методом, близким к круговому методу Харди - Литлвуда - Рамануджана (этот метод наз. круговы.м методом Харди - Литлвуда - $\mathrm{Pa}-$ мануджана в форме тригонометрич. сумм Виноградова). В большинстве других задач (распределение дробных долей, целые точки в областях и др.) главный член получается тривиально. Теперь возникает проблема оценки остаточного члена, и если удается доказать, что он является величиной меньшего порядка, чем главный, то тем самым и доказывается асимптотич. формула.



Основной задатей при оценке остаточного члена является задача вовможно более точных оценок тригонометрич. сумм. О методах оценок тригонометрич. сумм и оценках см. Тригонометрическая сумма, а также Виноградова метод, Варинга проблема, Гольд баха проблема, Аддитивиье проблемь.



Лит. см. при ст. Тригонометрическал сумма.



ТРИГОНОМЕТРИЧЕСКО $A$ A Kарачуба. НИЕ - приближенное представление функции $f(x)$ в виде тригонометрич. полинома



\[

T(x)=A+\sum_{k=1}^{n}\left(a_{k} \cos k x+b_{k} \sin k x\right),

\]



значения к-рого' в заданных точках совнадают с соответствующими значениями функции. Именно, всегда можно подобрать $2 n+1$ коәффициентов $A, a_{k}, b_{k}$, $k=0,1, \ldots, n$, полинома $n$-го порядка $T(x)$ так, ттобы его значения были равны зчачениям $y_{k}$ функции $y=f(x)$ в $2 n+1$ наперед заданных точках $x_{k}$ промежутка $[0,2 \pi)$. Полином $T(x)$ имест вид



\[

T(x)=\sum_{k=0}^{2 n} y_{k} t_{k}(x)

\]



где



$t_{k}(x)=\frac{\Delta x}{\Delta^{\prime}(x) 2 \sin \left(x-x_{k}\right) / 2}, \Delta(x)=\prod_{k=0}^{2 n} 2 \sin \frac{x-x_{k}}{2}$.



Особенно простой вид полином $T(x)$ приобретает в случае равноотстоящих узлов $x_{k}=(2 k-1) \pi / 2 n$; сго коэффициенты выражены формулами



\[

\begin{gathered}

A=-\frac{1}{2 n+1} \sum_{k=0}^{2 n} y_{k}, \\

a_{m}=\frac{2}{2 n+1} \sum_{k=0}^{2 n} y_{k} \cos m x_{k}, \\

b_{m}=\frac{2}{2 n+1} \sum_{k=0}^{2 n} y_{k} \sin m x_{k}, 1 \leqslant m \leqslant n .

\end{gathered}

\]



B. 1.. Битючков.



ТРИГОНОМЕТРИЯ - раздел геометрии, в к-ром метрич. соотношения между әлемештами треугольника описываются через тригонометрические бункции (см., напр., Косинусов теорема, Синусов теорема, Тангенсов формула), а также устанавливаются соотношения между тригонометрич. функциями (т. н. г о н и о ме тр пя). Т. рассматривается как для евклидовой, так и для неевклидовой геометрии. Т. сферы евклидового пространства наз. сферической тригонометрией.



ТРИКоМИ ЗАДАЧА - задача отыскангя решения уравнения смешанного эллиитико-гишерболического тиша с двумя независимыми переменными и с одной





\end{document}